\subsection{Extremal points of compact convex sets}

DEFINITION 1. --- Let A be a convex set in an affine space E. Then we say that a point \(x \in A\) is an extremal point of A if there does not exist an open segment that is contained in A and contains x.

In other words, the relations \(x=\lambda y+(1-\lambda)z\) , \(y\in A\) , \(z\in A\) , \(y\neq z\) and \(0\leqslant\lambda\leqslant1\) imply \(\lambda=0\) or \(\lambda=1\) (thus x=y or x=z). This implies that x cannot be the barycentre of a set of n points \(x_{i}\) of A carrying positive masses unless x is one of the \(x_{i}\) ; for this is just the definition when n=2; for arbitrary n argue by induction on n, as x is the barycentre of \(x_{1}\) and of the barycentre \(y_{1}\) of the \(x_{i}\) with \(2\leqslant i\leqslant n\) , therefore x is identical with \(x_{1}\) or \(y_{1}\) , and in the second case it is sufficient to apply the induction hypothesis.

To say that \(x\) is an extremal point of \(A\) also means that \(A - \{x\}\) is convex.

Examples. --- 1) In the space \(\mathbf{R}^n\) , all the points of the sphere \(\mathbf{S}_{n-1}\) are extremal points of the closed ball \(\mathbf{B}_n\) . For, if \(\sum_{i} y_i^2 \leqslant 1\) , \(\sum_{i} z_i^2 \leqslant 1\) and \(0 < \lambda < 1\) , the relation

\[
\lambda^ {2} \sum_ {i} y _ {i} ^ {2} + (1 - \lambda) ^ {2} \sum_ {i} z _ {i} ^ {2} + 2 \lambda (1 - \lambda) \sum_ {i} y _ {i} z _ {i} = 1 = (\lambda + (1 - \lambda)) ^ {2}
\]

is possible only if

\[
\sum_ {i} y _ {i} ^ {2} = \sum_ {i} z _ {i} ^ {2} = \sum_ {i} y _ {i} z _ {i} = 1.
\]

But this implies \(\sum_{i}(y_i - z_i)^2 = 0\) , thus \(y_i = z_i\) for all \(i\) , which proves our assertion.

\begin{enumerate}
\def\labelenumi{\arabic{enumi})}
\setcounter{enumi}{1}
\tightlist
\item
  In the normed space \(\mathcal{B}(\mathbf{N})\) of bounded sequences of real numbers (I, p.~4) the extremal points of the unit ball are the points \(x = (\xi_n)\) such that \(|\xi_n| = 1\) for all \(n\) . For, suppose that we had \(|\xi_n| \leqslant 1\) for all \(n\) and \(|\xi_p| < 1\) for one index \(p\) . We can then write
\end{enumerate}

\[
x = \frac {1 + \xi_ {p}}{2} y + \frac {1 - \xi_ {p}}{2} z
\]

where \(y\) (resp. \(z\) ) is the point all of whose coordinates are equal to the coordinate of \(x\) with the same index, except in the case of index \(p\) where the coordinate is equal to 1 (resp. - 1). This shows that \(x\) is not extremal, since we have \(\| y \| \leqslant 1\) and \(\| z \| \leqslant 1\) . Conversely, if \(|\xi_n| = 1\) for all \(n\) , then \(x\) is extremal, for the relation \(\xi_n = \lambda \eta_n + (1 - \lambda) \zeta_n\) with \(|\eta_n| \leqslant 1\) , \(|\zeta_n| \leqslant 1\) and \(0 < \lambda < 1\) implies \(\xi_n = \eta_n = \zeta_n\) .

\begin{enumerate}
\def\labelenumi{\arabic{enumi})}
\setcounter{enumi}{2}
\tightlist
\item
  Let \(u: \mathrm{E} \to \mathrm{E}'\) be an affine mapping of an affine space \(\mathrm{E}\) in an affine space \(\mathrm{E}'\) ; let \(\mathrm{C} \subset \mathrm{E}\) , \(\mathrm{C}' \subset \mathrm{E}'\) be two convex sets such that \(u(\mathrm{C}) \subset \mathrm{C}'\) . If \(x'\) is an extremal point of \(\mathrm{C}'\) and \(x\) is an extremal point of \(u^{-1}(x') \cap \mathrm{C}\) , then \(x\) is an extremal point of \(\mathrm{C}\) , as it follows from def. 1.
\end{enumerate}

PROPOSITION 1. --- Let B be the set of extremal points of A, a non-empty compact convex set in a Hausdorff locally convex space E, and let f be a convex function defined in A and upper semi-continuous. Then f attains its upper bound in A at one point (at least) of B.

Use F to denote the family of subsets X of A that are non-empty, closed, and such that every open segment that is contained in A and meets X necessarily lies in X. It has the following properties;

\begin{enumerate}
\def\labelenumi{(\roman{enumi})}
\item
  A belongs to \(\mathfrak{F}\) .
\item
  A point \(a \in \mathbf{A}\) is such that \(\{a\} \in \mathfrak{F}\) , if, and only if, \(a\) is an extremal point of \(\mathbf{A}\) .
\item
  Every non-empty intersection \(\mathbf{X}\) of a family \((\mathbf{X}_{\alpha})\) of sets of \(\mathfrak{F}\) also belongs to \(\mathfrak{F}\) .
\end{enumerate}

The properties (i), (ii) and (iii) follow immediately from the definitions.

\begin{enumerate}
\def\labelenumi{(\roman{enumi})}
\setcounter{enumi}{3}
\tightlist
\item
  Let \(X \in F\) , and let h be a function that is convex and upper semi-continuous in A; then the set Y of the points of X where the restriction \(h|X\) attains its upper bound in X is such that Y belongs to F.
\end{enumerate}

For, \(h|X\) being upper semi-continuous in \(X\) attains its upper bound \(\alpha\) over \(X\) in at least one point of \(X\) (GT, IV, § 6.2, th. 3); thus \(Y\) is non-empty, it is also closed (GT, IV, § 6.2, prop. 1). On the other hand let \(x, y\) be two distinct points of \(A\) and let \(z = \lambda x + (1 - \lambda)y\) be a point of \(Y\) such that \(0 < \lambda < 1\) ; as \(Y \subset X\) and \(X \in \mathfrak{F}\) , we have \(x \in X\) and \(y \in X\) ; on the other hand, as \(h\) is convex, we have

\[
h (z) \leqslant \lambda h (x) + (1 - \lambda) h (y)
\]

but as \(h(x) \leqslant \alpha\) , \(h(y) \leqslant \alpha\) and \(h(z) = \alpha\) , of necessity \(h(x) = h(y) = \alpha\) , that is to say \(x \in \mathbf{Y}\) and \(y \in \mathbf{Y}\) . Therefore \(\mathbf{Y} \in \mathfrak{F}\) .

With these properties established, let M be the set of \(x \in A\) where f attains its upper bound in A; by (iv), \(M \in F\) . On the other hand, by (iii) and the fact that the sets of F are closed subsets of the compact set A, it follows that F is inductive for the order relation \(\supset\) . By th. 2 of S, III, §2.4, M contains a subset N which is a minimal element of F. We shall show that N consists of a single point and this will complete the proof of the proposition. Since E is a Hausdorff locally convex space, it is sufficient to show that every continuous linear form u on E is constant in N (II, p.~38, cor. 1). Now it follows from (iv) that the set \(N'\) of the \(x \in N\) where \(u|N\) attains its upper bound in N is such that \(N'\) belongs to F; since N is minimal in F we necessarily have \(N' = N\) .

COROLLARY. --- Let A be a compact convex set in a Hausdorff locally convex space E. Then every closed support hyperplane H of A contains at least one extremal point of A.

For, if \(f(x) = \alpha\) is an equation of \(H\) and \(f(x) \leqslant \alpha\) in \(A\) , it is sufficient to apply prop. 1 to \(f\) .

THEOREM 1 (Krein-Milman). --- In a Hausdorff locally convex space E, every compact convex set A is the closed convex envelope of the set of its extremal points.

For, let C be the closed convex envelope of the set of extremal points of A; clearly \(C \subset A\) . To see that \(A \subset C\) , it is sufficient to prove that, if u is an affine linear function, continuous in E and if \(u(x) \geqslant 0\) in C then also \(u(x) \geqslant 0\) in A (II, p.~39, cor. 4); but this follows from prop. 1 applied to -u.

PROPOSITION 2. --- Let x be an extremal point of a compact convex set A in a Hausdorff locally convex space E. Then for every open neighbourhood V of x in E, there exists an open half-space F in E such that \(x \in F \cap A \subset V \cap A\) (in other words, the traces on A of the open half-spaces containing x, form a fundamental system of neighbourhoods of x in A).

For every open half-space D of E containing x, the set \(A \cap \overline{D}\) is a compact neighbourhood of x in A, and the intersection of all these neighbourhoods is precisely the point x (any two distinct points can be strictly separated by a closed hyperplane (II, p.~38, prop. 4). By prop. 1 of GT, I, § 9.2, it is sufficient to prove that the sets \(A \cap \overline{D}\) form a filter base. Now if we write \(L_{D} = A \cap (E - D)\) , the set \(L_{D}\) is convex, compact and contained in the convex set \(A - \{x\}\) ; if \(D_{1}, D_{2}\) are two open half-spaces of E containing x, the convex envelope B of \(L_{D_{1}} \cup L_{D_{2}}\) is therefore contained in \(A - \{x\}\) ; but B is a compact set (II, p.~14, prop. 15), therefore there exists a closed hyperplane H that separates x strictly from B (II, p.~38, prop. 4) and if the open half-space determined by H and containing x is D, then we have \(L_{D_{1}} \cup L_{D_{2}} \subset L_{D}\) , therefore \(A \cap \overline{D} \subset (A \cap \overline{D}_{1}) \cap (A \cap \overline{D}_{2})\) .

COROLLARY. --- In a Hausdorff locally convex space let K be a compact subset of a compact convex set A. Then the following conditions are equivalent.

\begin{enumerate}
\def\labelenumi{\alph{enumi})}
\item
  A is the closed convex envelope of K.
\item
  K meets every set that is the intersection of A with one of its support hyperplanes.
\item
  K contains the set of extremal points of A.
\end{enumerate}

\(a) \Rightarrow b)\) . Suppose that there exists a support hyperplane H of A whose equation is \(f(x) = \alpha\) , such that \((H \cap A) \cap K = \varnothing\) and suppose, for example, that \(f(x) \geqslant \alpha\) in A. As \(f(x) - \alpha > 0\) for all \(x \in K\) by hypothesis and as K is compact we have

\[
\beta = \inf _ {x \in \mathbf {K}} f (x) > \alpha ,
\]

and K is, therefore, contained in the closed half-space \(f(x) \geqslant \beta\) ; therefore the same is true of the closed convex envelope A of K and this is absurd.

\(b) \Rightarrow c)\) . Suppose that an extremal point \(x\) of \(\mathbf{A}\) does not belong to \(\mathbf{K}\) ; there is a neighbourhood \(\mathbf{V}\) of \(x\) in \(\mathbf{E}\) such that \(\mathbf{V} \cap \mathbf{A} \cap \mathbf{K} = \varnothing\) . But by prop. 2, we can suppose that \(\mathbf{V}\) is an open half-space defined by a hyperplane \(\mathbf{H}\) with the equation \(f(z) = \alpha\) . If for example \(f(x) > \alpha\) , then for all \(y \in \mathbf{K}\) , we have \(f(y) \leqslant \alpha\) , therefore \(\mathbf{K}\) does not meet the intersection of \(\mathbf{A}\) and the support hyperplane \(f(z) = \gamma > \alpha\) parallel to \(\mathbf{H}\) (II, p.~37, prop. 2); this is absurd.

\(c) \Rightarrow a)\) . This is an obvious consequence of the Krein-Milman theorem.

exerc. 11). The proof of theorem 1 (II, p.~56) shows that in any case A is the convex closed envelope of the set of extremal points of A which belong to a support hyperplane.

